\documentclass[../main.tex]{subfiles}  % 使用主文件的文档类

\begin{document}

\section{Discussion and Conclusion}

This paper presented an extensive analysis of security vulnerabilities in wireless communication protocols, specifically focusing on Bluetooth Classic (BC), Bluetooth Low Energy (BLE), and WPA2. Through detailed case studies, we identified critical weaknesses in session establishment and key negotiation processes that could be exploited by adversaries to compromise the confidentiality, integrity, and availability of wireless communications.

Our findings underscore the urgent need for the adoption of stronger cryptographic standards and more rigorous security practices in protocol design and implementation. The continued reliance on outdated and insecure methods, driven by legacy support and usability concerns, represents a significant threat to the security landscape of wireless communications. To address these challenges, we recommend the following: the enforcement of robust encryption standards across all devices, the phasing out of deprecated security protocols, and the implementation of adaptive security mechanisms that can dynamically respond to emerging threats.

Moreover, the transition to more secure protocols like WPA3 should be prioritized, supported by industry-wide efforts to educate stakeholders on the importance of updating and securing wireless infrastructure. As the ecosystem of connected devices continues to expand, the implications of insecure wireless communications will only grow more severe, necessitating a proactive approach to security by design.

In conclusion, enhancing the security of wireless communication protocols is a multifaceted challenge that requires coordinated efforts across academia, industry, and regulatory bodies. By advancing our understanding of existing vulnerabilities and developing innovative solutions, we can better protect the integrity of wireless communications in an increasingly interconnected world.



\end{document}
